\begin{abstract}
Perhitungan PPh Pasal 21 pegawai tetap berubah hampir setiap tahun. Sejak 2024 metode disetahunkan digantikan tarif efektif rata-rata (TER) bulanan, tahun 2025--2026 menambah fasilitas PPh ditanggung pemerintah (DTP), dan parameter BPJS diperbarui setiap tahun. Di saat yang sama, setiap perusahaan memiliki kebijakan payroll sendiri. Dalam praktik, kedua jenis aturan itu ditanam bercampur di rumus \en{spreadsheet} atau kode program. Akibatnya perubahan regulasi menuntut perubahan rumus, kebijakan perusahaan dapat melanggar aturan wajib tanpa terdeteksi, dan angka di slip gaji tidak dapat ditelusuri ke pasal yang mendasarinya. Penelitian ini mengusulkan \en{knowledge-based system} untuk perhitungan PPh Pasal 21 tahun pajak 2023--2026 dengan tiga gagasan utama. Pertama, \en{knowledge base} deklaratif dua lapis, yaitu lapisan regulasi dan lapisan kebijakan perusahaan, dengan resolusi konflik berbasis asas \en{lex superior}, \en{lex specialis}, dan \en{lex posterior}. Kedua, versi waktu bitemporal (\en{valid time} dan \en{knowledge time}), sehingga satu \en{inference engine} \en{forward chaining} yang sama melayani dua rezim pajak dan regulasi yang berlaku surut. Ketiga, perlakuan deklaratif atas definisi melingkar \en{gross-up} sebagai \en{fixed point} dengan jaminan Knaster--Tarski. Setiap angka disertai \en{explanation trace} hingga aturan dan pasal sumbernya. Sistem dilengkapi demo Streamlit dan aplikasi web Laravel untuk admin \en{finance} yang memanggil \en{engine} lewat jembatan JSON tanpa menyalin rumus pajak. Karena penelitian ini tidak melibatkan \en{domain expert} maupun data payroll riil, kebenaran ditetapkan lewat \en{oracle pyramid}: 45 contoh resmi regulasi dengan 1.213 nilai harapan dan protokol erratum, \en{differential testing} terhadap \en{baseline} \en{hard-coded} yang dibekukan sebelum aturan \en{knowledge base} ditulis, \en{metamorphic relation} yang sebagian besar diuji dengan \en{property-based testing}, dan \en{mutation testing} atas \en{knowledge base}. Seluruh perhitungan tunduk pada kontrak presisi tanpa \en{floating point}. Draf ini menyajikan rancangan sistem dan metodologi evaluasi dua belas eksperimen. Hasil eksperimen belum final dan akan dilaporkan terpisah.
\end{abstract}

\begin{IEEEkeywords}
\en{knowledge-based system}, \en{rules as code}, PPh Pasal 21, \en{forward chaining}, \en{bitemporal versioning}, \en{conflict resolution}, \en{metamorphic testing}, \en{explainability}
\end{IEEEkeywords}
